\section{Project Planning \& Task Distribution}

\subsection{Kế hoạch triển khai dự án}

Dự án Bloxorz Solver được chia thành các nhóm công việc chính: xây dựng
logic trò chơi, phát triển giao diện, cài đặt các thuật toán tìm kiếm và
chuẩn bị báo cáo. Logic trò chơi và mô hình trạng thái là nền tảng dùng
chung, giúp các thành viên phát triển giao diện và thuật toán song song
mà vẫn sử dụng thống nhất các quy tắc di chuyển, công tắc và cầu.

Kế hoạch triển khai dưới đây được tổng hợp từ lịch sử commit trong
khoảng ngày 12/10/2026 đến 22/10/2026. Các giai đoạn có thời gian giao
nhau do nhiều mô-đun được phát triển đồng thời.

\begin{table}[H]
    \centering
    \caption{Các giai đoạn triển khai dự án}
    \label{tab:project_timeline}
    \small
    \renewcommand{\arraystretch}{1.2}
    \begin{tabularx}{\textwidth}{@{}>{\raggedright\arraybackslash}p{2.7cm}
        >{\raggedright\arraybackslash}p{2.5cm}
        >{\raggedright\arraybackslash}X@{}}
        \toprule
        \textbf{Giai đoạn} & \textbf{Thời gian} & \textbf{Nội dung chính} \\
        \midrule
        Xây dựng nền tảng & 12--15/10 &
        Khởi tạo repository và cấu trúc mã nguồn; xây dựng các đối tượng
        trò chơi, \texttt{Board}, \texttt{GameRules}, dữ liệu 20 màn chơi
        và các kiểm thử logic. \\
        \midrule
        Phát triển giao diện & 17--19/10 &
        Xây dựng ứng dụng Pygame, bộ điều khiển và hiển thị bàn chơi;
        bổ sung menu, quản lý tiến trình, hiệu ứng chuyển động và sửa
        lỗi tương tác. \\
        \midrule
        Phát triển thuật toán & 19--22/10 &
        Mô hình hóa bài toán; xây dựng khung solver dùng chung; cài đặt
        BFS, DFS, UCS và A*; bổ sung heuristic, mã đo hiệu năng và
        kiểm thử thuật toán. \\
        \midrule
        Viết báo cáo& Từ ngày 19 &
        Viết báo cáo LaTeX, chuẩn bị video\\
        \bottomrule
    \end{tabularx}
\end{table}

\subsection{Phân công nhiệm vụ}

\begin{table}[H]
    \centering
    \caption{Phân công nhiệm vụ của các thành viên}
    \label{tab:task_distribution}
    \small
    \renewcommand{\arraystretch}{1.2}
    \begin{tabularx}{\textwidth}{@{}>{\raggedright\arraybackslash}p{3.5cm}
        >{\raggedright\arraybackslash}X@{}}
        \toprule
        \textbf{Thành viên} & \textbf{Nhiệm vụ} \\
        \midrule
        Ngô Hoàng Minh &
        \textbullet\ Khởi tạo dự án, tổ chức mã nguồn và dữ liệu dùng chung. \newline
        \textbullet\ Xây dựng \texttt{Board}, \texttt{GameRules}, bộ 20 màn chơi và công cụ trực quan hóa bản đồ. \newline
        \textbullet\ Mô hình hóa bài toán, xây dựng khung solver, cài đặt A* và heuristic. \newline
        \textbullet\ Viết kiểm thử logic, bản đồ và A*; sửa lỗi GUI. \newline \\
        \midrule
        Phan Nhật Minh &
        \textbullet\ Định nghĩa các loại ô và cầu; khai báo thư viện phụ thuộc. \newline
        \textbullet\ Cài đặt thuật toán BFS. \newline
        \textbullet\ Bổ sung mã đo thời gian, bộ nhớ, số nút mở rộng và độ dài lời giải. \\
        \midrule
        Nguyễn Thành Bảo &
        \textbullet\ Xây dựng cơ chế khối, công tắc và tương tác với cầu. \newline
        \textbullet\ Tích hợp các thành phần vào API backend dùng chung. \newline
        \textbullet\ Cài đặt UCS, viết kiểm thử và tài liệu thuật toán. \\
        \midrule
        Đỗ Ngọc Hải &
        \textbullet\ Xây dựng ứng dụng Pygame, bộ điều khiển và hiển thị trò chơi. \newline
        \textbullet\ Quản lý tiến trình, sửa hiệu ứng rơi và viết kiểm thử giao diện, chuyển động. \newline
        \textbullet\ Cài đặt DFS, bổ sung cấu trúc kết quả tìm kiếm và kiểm thử DFS. \\
        \midrule
        Nguyễn Xuân Lộc &
        \textbullet\ Điều chỉnh tỷ lệ hiển thị bàn chơi, khối và ký hiệu công tắc. \newline
        \textbullet\ Hoàn thiện menu, màn hình chiến thắng và hiệu ứng ngôi sao. \newline
        \textbullet\ Cải thiện hiệu ứng rơi và bổ sung cơ chế tự khởi động lại màn chơi khi thua. \\
        \bottomrule
    \end{tabularx}
\end{table}

